\documentclass[10pt]{beamer}
%%%%%%%%%%%%%%%%%%%%%%%%%%%%%%%%%%%
% File: preamble.tex
%%%%%%%%%%%%%%%%%%%%%%%%%%%%%%%%%%%

\usepackage{geometry}
\geometry{left = 1cm, right = 1cm, top = 1.5cm, bottom = 1.5cm}

% Set fonts commands
\newcommand{\song}{\CJKfamily{song}} 
\newcommand{\hei}{\CJKfamily{hei}} 
\newcommand{\kai}{\CJKfamily{kai}} 
\newcommand{\fs}{\CJKfamily{fs}}
\newcommand{\tqs}{\textquotesingle}
\usepackage{fontawesome}

\newcommand{\me}[2]{\author{{\bfseries 姓名：}\underline{#1}\hspace{2em}{\bfseries 学号：}\underline{#2}}}

% Always keep this.
\newcommand{\noplagiarism}{
  \begin{center}
    \fbox{\begin{tabular}{@{}c@{}}
      请独立完成作业，不得抄袭。\\
      若参考了其它资料，请给出引用。\\
      鼓励讨论，但需独立书写解题过程。
    \end{tabular}}
  \end{center}
}

% Each hw consists of three parts:
% (1) this homework
\newcommand{\beginthishw}{\part{作业}}
% (2) corrections (Optional)
\newcommand{\begincorrection}{\part{订正}}
% (3) any feedback (Optional)
\newcommand{\beginfb}{\part{反馈}}

% For math
\usepackage{amsmath}
\usepackage{amsfonts}
\usepackage{amssymb}
\usepackage{graphicx}
\usepackage{listings}
\usepackage{xcolor}
%\usepackage{clrscode3e} 
\usepackage{enumitem}
\usepackage{tikz}
\usepackage{tabularx}
\usepackage{multirow}
\newcolumntype{Y}{>{\centering\arraybackslash}X}
\newcolumntype{P}{>{\centering\arraybackslash}p}
%bigger integrate symbol
\usepackage{exscale}
\usepackage{relsize}
\usepackage{textcomp}

% colors
\newcommand{\red}[1]{\textcolor{red}{#1}}
\newcommand{\blue}[1]{\textcolor{blue}{#1}}
\newcommand{\teal}[1]{\textcolor{teal}{#1}}

% algorithms
\usepackage[]{algorithm}
\usepackage[noend]{algpseudocode} % noend
\makeatletter
\newenvironment{breakablealgorithm}
  {% \begin{breakablealgorithm}
   \begin{center}
     \refstepcounter{algorithm}% New algorithm
     \hrule height.8pt depth0pt \kern2pt% \@fs@pre for \@fs@ruled 画线
     \renewcommand{\caption}[2][\relax]{% Make a new \caption
       {\raggedright\textbf{\ALG@name~\thealgorithm} ##2\par}%
       \ifx\relax##1\relax % #1 is \relax
         \addcontentsline{loa}{algorithm}{\protect\numberline{\thealgorithm}##2}%
       \else % #1 is not \relax
         \addcontentsline{loa}{algorithm}{\protect\numberline{\thealgorithm}##1}%
       \fi
       \kern2pt\hrule\kern2pt
     }
  }{% \end{breakablealgorithm}
     \kern2pt\hrule\relax% \@fs@post for \@fs@ruled 画线
   \end{center}
  }
\makeatother
\renewcommand{\algorithmicrequire}{\textbf{Input:}} % Use Input in the format of Algorithm
\renewcommand{\algorithmicensure}{\textbf{Output:}} % Use Output in the format of Algorithm
% See [Adjust the indentation whithin the algorithmicx-package when a line is broken](https://tex.stackexchange.com/a/68540/23098)
\newcommand{\algparbox}[1]{\parbox[t]{\dimexpr\linewidth-\algorithmicindent}{#1\strut}}
\newcommand{\hStatex}[0]{\vspace{5pt}}
\makeatletter
\newlength{\trianglerightwidth}
\settowidth{\trianglerightwidth}{$\triangleright$~}
\algnewcommand{\LineComment}[1]{\Statex \hskip\ALG@thistlm \(\triangleright\) #1}
\algnewcommand{\LineCommentCont}[1]{\Statex \hskip\ALG@thistlm%
  \parbox[t]{\dimexpr\linewidth-\ALG@thistlm}{\hangindent=\trianglerightwidth \hangafter=1 \strut$\triangleright$ #1\strut}}
\makeatother


% Define theorem-like environments
\usepackage[amsmath, thmmarks, framed]{ntheorem}
\usepackage{framed}

\theoremheaderfont{\kai\bfseries}
\theoremstyle{break}
\theorembodyfont{\song}
\theoremseparator{\vspace{1mm}}
\renewcommand*\FrameCommand{{\color{gray}\vrule width 3pt \hspace{10pt}}}
\newframedtheorem{problem}{\faCheckSquareO \ Problem}

\theorempostwork{\hrule}
\newtheorem*{solution}{\faEdit \ Solution}
\newtheorem*{revision}{\faEdit \ Revision}

\theoremstyle{plain}
\newtheorem*{cause}{\faCoffee \ Cause}
\newtheorem*{remark}{\faCommentingO \ Remark}

\theoremstyle{break}
\theoremsymbol{\ensuremath{\Box}}
\newtheorem*{proof}{\faEdit \ Proof}

\renewcommand\figurename{Figure}
\renewcommand\tablename{Table}

%enumeration
\setenumerate[1]{
	itemsep=0pt,
	partopsep=0pt,
	parsep=\parskip,
	topsep=0pt,
	leftmargin=20pt
}
\setitemize[1]{
	itemsep=0pt,
	partopsep=0pt,
	parsep=\parskip,
	topsep=0pt,
	leftmargin=20pt
}
\setdescription{
	itemsep=0pt,
	partopsep=0pt,
	parsep=\parskip,
	topsep=0pt,
	leftmargin=20pt
}
\lstset{
	language={[ISO]C++},
    numbers=left, 
    numberstyle= \tiny,
    commentstyle=\color{red!50!green!50!blue!50},  
  	rulesepcolor=\color{red!20!green!20!blue!20},  
 	keywordstyle=\color{blue!90}\bfseries,   
  	showstringspaces=false, 
  	stringstyle=\ttfamily,
} 


% For figures
% for fig with caption: #1: width/size; #2: fig file; #3: fig caption
\newcommand{\fig}[3]{
	\centerline{\includegraphics[scale=#1]{#2}}
	\centerline{#3}
}

% for fig without caption: #1: width/size; #2: fig file
\newcommand{\fignc}[2]{
	\centerline{\includegraphics[scale=#1]{#2}}
}

%%%%%%%%%%%%%%%%%%%%%%%%%%%%%%%%%%%%%%%%%%%%%%%%%%%%%%%%%%%

%title
\title{Beamer Metropolis Mod README}
\author[Tianyun Zhang]{张天昀}
\institute[CS@NJU]{南京大学计算机科学与技术系\\ 171860508@smail.nju.edu.cn}
\renewcommand{\today}{\number\year 年\number\month 月\number\day 日}
\date{\today}

\renewcommand{\figurename}{图}
\renewcommand{\tablename}{表}


\begin{document}

\begin{frame}
  \maketitle
\end{frame}

\begin{frame}{说明}
    由于NJU beamer模板实在太丑了，于是我找了这个 3000+ Star 的 \href{https://github.com/matze/mtheme}{beamer 主题}来当我的新的beamer模板，顺便写一下beamer很多东西用法的memo。
    
    \begin{block}{LICENSE}
        本文档使用CC-BY-SA 4.0协议授权。
        \ccbysa
    \end{block}
\end{frame}

\begin{frame}{目录}
    \setbeamertemplate{section in toc}[sections numbered]
    \tableofcontents[hideallsubsections]
\end{frame}

\section{基础内容}
\begin{frame}[fragile]{页面布局}
    一般beamer总是从左往右、从上至下（居中）来分布页面内容，但是可以手动分栏。
    \begin{verbatim}
        \begin{columns}[T,onlytextwidth]
            \column{0.5\textwidth}
                Foo
            \column{0.5\textwidth}
                Bar
        \end{columns}
    \end{verbatim}
    就会产生如下效果：
    \begin{columns}[T,onlytextwidth]
        \column{0.5\textwidth}
        Foo
        \column{0.5\textwidth}
        Bar
    \end{columns}
\end{frame} 

\begin{frame}[fragile]{页面元素}
    页面元素有很多种，除了文字以外，有
    公式：$$y = x^2$$
    图片：\vspace{2mm} \\
    \begin{center}
        \includegraphics[scale=0.1]{demo.jpg}
    \end{center}
    
    表格：\vspace{2mm} \\
    \begin{center}
        \begin{tabular}{|c|c|c|}
            \hline
            \multicolumn{2}{|c|}{111} & \multirow{2}{*}{FooBar} \\
            \cline{1-2}
            Foo & Bar &  \\
            \hline
        \end{tabular}
    \end{center}
\end{frame} 

\begin{frame}{页面元素}
    tikz绘制的图表：
    \begin{figure}
        \begin{tikzpicture}
            \begin{axis}[
                mbarplot,
                xlabel={Foo},
                ylabel={Bar},
                width=0.9\textwidth,
                height=6cm,
            ]

            \addplot plot coordinates {(1, 20) (2, 25) (3, 22.4) (4, 12.4)};
            \addplot plot coordinates {(1, 18) (2, 24) (3, 23.5) (4, 13.2)};
            \addplot plot coordinates {(1, 10) (2, 19) (3, 25) (4, 15.2)};

            \legend{lorem, ipsum, dolor}

            \end{axis}
        \end{tikzpicture}
    \end{figure}
\end{frame}

\begin{frame}[fragile]{页面元素}
    代码块（页面需要fragile）：
    \begin{verbatim}
        int main(){return 0;}
    \end{verbatim}
    引言：
    \begin{quote}
        同*恋才是真爱。——鲁迅
    \end{quote}
    \begin{quote}
        ‘The Kouga Ninja Scrolls’ is a historical fantasy novel about ninja written in 1958-1959 by the Japanese author Futaro Yamada. This is the first volume of the Ninja Scrolls series written by Yamada in 1958-2001. The book has been translated into English by Geoff Sant and was published by Del Rey in December 2006.
        \begin{flushright}
            -- from Wikipedia, the free encyclopedia
        \end{flushright}
    \end{quote}
\end{frame}

\begin{frame}[standout]
    \textsl{The end.}
\end{frame}
\end{document}